\documentclass[10pt]{article}

% , , , - fonts & encoding , , , -
\usepackage[T1]{fontenc}
\usepackage[utf8]{inputenc}

\usepackage{mathptmx}          % Times text + math (ICLR/NeurIPS look); tt stays Latin Modern (scalable)

% , , , - layout , , , -
\usepackage[letterpaper,margin=1in]{geometry}
\usepackage[protrusion=true,expansion=false]{microtype}
\usepackage[round]{natbib}
\bibliographystyle{plainnat}

% , , , - math & tables , , , -
\usepackage{amsmath,amssymb,mathtools}
\usepackage{booktabs,array,multirow}
\usepackage{graphicx}
\usepackage{caption}
\usepackage{subcaption}
\usepackage{enumitem}
\usepackage{xcolor}
\usepackage{tikz}
\usetikzlibrary{arrows.meta,positioning,shapes.geometric,calc}
\usepackage{listings}
\usepackage[colorlinks=true,linkcolor=blue!45!black,citecolor=blue!45!black,urlcolor=blue!45!black]{hyperref}

% , , , - title spacing (compact, conference-style) , , , -
\usepackage{titlesec}
\titlespacing*{\section}{0pt}{1.05\baselineskip}{0.35\baselineskip}
\titlespacing*{\subsection}{0pt}{0.7\baselineskip}{0.25\baselineskip}
\titlespacing*{\paragraph}{0pt}{0.4\baselineskip}{0.8em}
\captionsetup{font=small,labelfont=bf}

% , , , - code listing style , , , -
\lstset{
  basicstyle=\ttfamily\small,
  keywordstyle=\color{blue!55!black},
  commentstyle=\color{green!40!black}\itshape,
  language=Python,
  showstringspaces=false,
  columns=fullflexible,
  keepspaces=true,
  frame=none,
  xleftmargin=1.2em
}

% , , , - macros , , , -
\newcommand{\hb}{h_{\ell}}
\newcommand{\Rmat}{\mathbf{R}}
\newcommand{\Real}{\mathbb{R}}
\newcommand{\Bind}{B}
\newcommand{\acc}{\mathrm{Acc}}
\newcommand{\cm}{\textsf{context\_matched}}
\newcommand{\says}{\textsf{says\_installed}}
\newcommand{\code}[1]{\texttt{#1}}
\emergencystretch=3em \tolerance=1500

\title{\vspace{-3.5em}\bfseries Security-Grounded Reasoning in Code Language Models:\\[0.15em]
From Decodable Structure to Causal Use}

\author{Anonymous authors\\\small Paper under double-blind review}
\date{}

\begin{document}
\maketitle
\vspace{-2.2em}

\begin{center}
\fcolorbox{red!60!black}{red!5}{\parbox{0.92\linewidth}{\textbf{Archived manuscript.}
The lens results in this version use superseded prompt-specific cotangent and
conserving-cotangent methods, not the published J-lens and R-lens. The active
paper and main results are maintained in \texttt{docs/PAPER\_DRAFT.md} and
\texttt{docs/RESULTS.md}.}}
\end{center}

\begin{abstract}
\noindent
Code models increasingly inspect vulnerabilities, modify repositories, and
operate development tools. Trust in these systems requires evidence that their
answers depend on program semantics when lexical cues are misleading. We study
variable binding, the relation between a variable use and its active definition,
in DeepSeek-Coder 1.3B and 6.7B and StarCoder2 3B. Paired programs differ by one
earlier character that changes the binding while preserving the queried token,
its position, and its local context. A bounded surface reader therefore has
exactly $0.500$ accuracy. We measure linear recoverability, transfer across
program transformations, causal use, attribution to the active definition, and
alignment with output vocabulary. Binding reaches $0.984$ linear-probe accuracy
in middle layers. Frozen probes transfer across irrelevant context and consistent
renaming, with lower accuracy under competing scopes and control-flow flattening.
We then learn a one-dimensional intervention that replaces the host program's
binding component with the component from a counterfactual program. The edited
model returns the value selected by the donor binding on every held-out example
for DeepSeek-Coder 6.7B and StarCoder2 3B. Reversing the value assignment between
training and evaluation rules out a fixed answer-token direction. A conservative
relevance decomposition assigns up to $22\%$ more of the answer score to the
definition that enters scope. A Jacobian-based vocabulary readout follows the
binding, although scope remains confounded with position and replacement in the
single template. These results provide bounded evidence for security-relevant
semantic reasoning. They also identify transformations under which that evidence
weakens.
\end{abstract}

% =====================================================================
\section{Introduction}
% =====================================================================
Code models now inspect vulnerabilities, edit repositories, review patches, and
operate build systems. Accuracy on familiar benchmarks provides limited evidence
for safe use in these settings. Identifier names, formatting, templates, nearby
tokens, and common control-flow shapes often correlate with program behavior.
These cues can fail on unfamiliar, obfuscated, or adversarial code. Reliable use
therefore requires evidence that model outputs depend on the program relations
that determine execution.

We study this requirement through one minimal relation, \emph{variable binding}:
which definition a variable use resolves to. Binding admits a construction that
isolates semantics from surface form exactly. In the pair below, the queried
token \code{x} is identical and occupies the same position, but one character in
an earlier assignment changes whether the use resolves to the outer or the inner
definition:
\begin{center}
\begin{minipage}{0.44\linewidth}
\begin{lstlisting}
x = 3
def f():
    y = 7
    return x   # returns 3
\end{lstlisting}
\end{minipage}\hfill
\begin{minipage}{0.44\linewidth}
\begin{lstlisting}
x = 3
def f():
    x = 7
    return x   # returns 7
\end{lstlisting}
\end{minipage}
\end{center}
A reader restricted to a bounded window around the queried token cannot
distinguish these cases. Predicting the value requires scope resolution and a
link from the use to the active definition. This construction supports separate
tests of representation, robustness, and causal use.

Prior analysis of code models tends to answer one of these questions in
isolation. Probing establishes that a fact is \emph{decodable}
\citep{karmakar2021whatdo,wan2022theyknow,troshin2022probing}; robustness studies
establish that behavior does or does not \emph{survive} a transformation
\citep{yefet2020adversarial,henkel2022semantic,rabin2021generalizability};
attribution and vocabulary lenses describe \emph{where} an output score comes from
\citep{bach2015lrp,nostalgebraist2020logitlens}. These are not
interchangeable. A probe can decode information the model never uses; a
behavioral task can be solved by a shortcut; an attribution map can look
plausible without being causal; and a causally important feature need not be
human-readable. We therefore treat semantic grounding as a conjunction of
separable claims and test each with an instrument matched to it.

We define grounding operationally. Let a program property $S(x)$ be fixed by
programming-language semantics (here, which definition a use resolves to). A
model is \emph{semantically grounded} with respect to $S$ to the extent that its
internal computation (i) is sensitive to changes in $S$ that leave surface form
fixed, (ii) is invariant to changes that preserve $S$, and (iii) causally uses an
internal representation of $S$ when producing the relevant output. This is
stronger than benchmark accuracy and decodability. We test five questions.
\textbf{(1) Representation:} Is binding linearly recoverable above a measured
surface floor? \textbf{(2) Robustness:} Does the frozen representation survive
meaning-preserving perturbations, and does structural interference reduce its
transfer? \textbf{(3) Causal use:} Does replacing a compact binding component
change the answer according to the installed binding? \textbf{(4) Attribution:}
Is the unedited answer score assigned to the active definition?
\textbf{(5) Lexical expression:} Does the representation align with scope terms
in output coordinates?

\paragraph{Findings and contributions.} Binding becomes linearly available in
early-to-middle layers, survives many surface changes, and is read causally from
a rank-1 component. Competing scope and flattened control flow reduce frozen-probe
transfer. The lexical alignment does not identify scope independently of position.
We contribute a paired construction with an exact $0.500$ surface floor, a
stratified probing protocol, a frozen-transfer analysis, a crossed $2\times2$
causal intervention, a conservative relevance decomposition, and a crossed-arm
vocabulary test. Together, these measurements support a bounded trust claim and
state where it fails.

% =====================================================================
\section{Related work}
% =====================================================================
\paragraph{Code models and what they encode.} Pretrained code models support
search, generation, and completion \citep{feng2020codebert,wang2021codet5,
chen2021codex,nijkamp2023codegen,li2023starcoder,lozhkov2024starcoder2,
guo2024deepseekcoder}, and one line explicitly injects data-flow structure at
pretraining time \citep{guo2021graphcodebert}. Probing studies report that such
models encode aspects of syntax and data flow
\citep{karmakar2021whatdo,wan2022theyknow,troshin2022probing,lopez2022astprobe},
and \citet{jin2023evidence} argue that semantic content emerges in models trained
only to generate programs. Recent work tests functional equivalence
\citep{maveli2025functional}, execution-trace information
\citep{wang2025codesemantics}, and internal signals for code correctness and
security \citep{tahimic2025correctness,vu2025autoprobe}. These studies strengthen
the case for white-box analysis of code models. Task and probe accuracy do not
establish causal use of program semantics. Our intervention addresses that gap.

\paragraph{Probing and its controls.} Linear probes read information already made
linearly available by the model \citep{alain2016understanding}, but probe
accuracy must be interpreted against controls for probe capacity and lexical
memorization \citep{hewitt2019control,pimentel2020information,voita2020information,
belinkov2022probing}. We adopt these controls, shuffled-label selectivity,
grouped splits, an embedding-layer reference, and add a model-free surface reader
evaluated on the same items, so that every claim is stated relative to a
\emph{measured} floor rather than a nominal chance level. Our
$0.500$ floor is not assumed; it is pinned by construction.

\paragraph{Causal and mechanistic interpretability.} Causal abstraction
distinguishes information that is present from information used with the role of
an interpretable variable \citep{geiger2021causal,vig2020causal}. Interventions
that localize or edit computation include activation patching and its analyses
\citep{meng2022rome,wang2023ioi,conmy2023acdc,zhang2024patching,heimersheim2024patching},
within the residual-stream view of \citet{elhage2021mathematical}. Distributed
Alignment Search (DAS) learns a subspace in which interchange interventions align
a distributed representation with a high-level causal variable
\citep{geiger2024das,wu2023boundless}. We apply DAS to the variable ``which
definition is in scope'' and add a crossed value assignment that forces the
correct output direction to reverse between fitting and evaluation, which
separates binding transport from an answer-token direction. The most closely
related targets, symbol binding and entity tracking in language models
\citep{smolensky1990tensor,feng2023entities,kim2023entity,prakash2024finetuning}%
, concern natural-language entities; we study the programming-language relation,
where ground truth is exactly decidable.

\paragraph{Vocabulary lenses and attribution.} The logit lens reads a hidden
state through the unembedding \citep{nostalgebraist2020logitlens}; the tuned lens
corrects it with a learned affine map \citep{belrose2023tunedlens}, and related
work reads intermediate states in vocabulary space
\citep{geva2022ffn,dar2023embedding}. The Jacobian lens averages downstream
Jacobians across positions and a broad prompt distribution to identify directions
that support later verbalization \citep{gurnee2026workspace}. The R-lens replaces
its raw-gradient backward pass with relevance-propagation coefficients
\citep{blank2026rlens}, following Relevance Patching
\citep{jafari2025relp}. Our vocabulary readout uses a prompt-specific downstream
Jacobian. It is related to the Jacobian lens and implements a different estimator.
Separately, layer-wise relevance
propagation \citep{bach2015lrp} has been adapted to transformers with
conservation-respecting rules for normalization and attention
\citep{ali2022xai,achtibat2024attnlrp}, in contrast to gradient attributions
\citep{sundararajan2017ig,chefer2021transformer}. We use both families as
\emph{descriptive} instruments: an archived prompt-specific cotangent vocabulary readout and
a conservative relevance decomposition. The latter predates and differs from the
2026 R-lens because it decomposes one answer score across source tokens. Neither
instrument establishes causal necessity. DAS supplies the causal test.

\paragraph{Robustness of code models.} Semantics-preserving transformations and
adversarial renamings degrade code-model behavior
\citep{yefet2020adversarial,bielik2020robust,henkel2022semantic,rabin2021generalizability}.
We connect this literature to representation: we freeze a probe trained on clean
programs and measure transfer under an execution-verified obfuscation ladder,
which distinguishes lexical fragility (renaming) from structural fragility
(control-flow flattening). Ground truth throughout is derived from a code property
graph \citep{yamaguchi2014cpg}.

% =====================================================================
\section{Problem formulation}
\label{sec:formulation}
% =====================================================================
Let $x$ be a program, $u$ the token position of a variable use, and
$\Bind(x,u)\in\{0,1\}$ the binding fixed by language semantics: $0$ for the outer
definition, $1$ for the inner. A frozen autoregressive model produces a hidden
state $\hb(x,u)\in\Real^{d}$ at layer $\ell$, read at the last token that covers
the use under causal attention. A paired generator constructs $x^{(0)}$ and
$x^{(1)}$ with $\Bind(x^{(0)},u)\neq \Bind(x^{(1)},u)$ such that the token at $u$,
its index, and its bounded local context are identical; the pair differs at one
earlier identifier character.

We call a property \emph{semantic} to the extent that a bounded reader of the
program's surface form cannot recover it. This definition is enforced, not
estimated: the paired construction makes the model-free reader score exactly
$0.500$, so any accuracy above that floor is computed by the model rather than
read off the text. The definition is falsifiable and has excluded a candidate
relation: control dependence has a \emph{measured} surface floor near $0.93$ (a
statement's guard is usually its nearest enclosing \code{if}), so decoding it at
ceiling constrains nothing about semantic representation, and we do not report it
(Appendix~\ref{app:controldep}). The floor is always stated relative to a named
reader (\S\ref{sec:setup}).

\paragraph{The five properties.} We operationalize grounding as five measurements.

\emph{(P1) Recoverability.} A linear probe $q_\ell$ predicts $\Bind$ from $\hb$
above the surface baseline $s$:
$\acc\!\big(q_\ell(\hb),\Bind\big)-\acc\!\big(s,\Bind\big) > 0.$

\emph{(P2) Transfer.} For a semantics-preserving transformation $T$, a probe
fitted on clean programs is frozen and re-evaluated,
$\Delta_T(\ell)=\acc\!\big(q_\ell(\hb(T(x),u_T)),\Bind\big)-\acc\!\big(q_\ell(\hb(x,u)),\Bind\big).$
Grounding requires both invariance to meaning-preserving $T$ and sensitivity to
meaning-changing $T$; indiscriminate invariance is as uninformative as
indiscriminate sensitivity.

\emph{(P3) Causal use.} DAS learns an orthonormal $\Rmat\in\Real^{d\times r}$ and
performs an interchange between a host and a donor state,
\begin{equation}
\hb' \;=\; \hb^{\mathrm{host}} + \Rmat\Rmat^{\!\top}\!\big(\hb^{\mathrm{donor}}-\hb^{\mathrm{host}}\big).
\label{eq:interchange}
\end{equation}
If $\Rmat$ realizes binding, the intervened model emits the value selected by the
donor binding, including when the map from binding to answer token is reversed in
a held-out arm.

\emph{(P4) Attribution.} For a selected answer score $f_y(x)$, a relevance
decomposition assigns token contributions $r_i$ satisfying approximate
conservation $\sum_i r_i \approx f_y(x)$; we ask whether changing $\Bind$
reallocates relevance from the definition becoming inactive to the one becoming
active, including over tokens that do not change.

\emph{(P5) Lexical expression.} For a predeclared token pair
$(w_{\mathrm{in}},w_{\mathrm{out}})$, a vocabulary lens $J_\ell$ gives a margin
$m_\ell(x)=J_\ell(\hb(x,u))_{w_{\mathrm{in}}}-J_\ell(\hb(x,u))_{w_{\mathrm{out}}}$;
the counterfactual shift $\delta=m_\ell(x^{(1)})-m_\ell(x^{(0)})$ is predicted to
be positive, and its sign is measured separately in the two value arms.

% =====================================================================
\section{Experimental setup}
\label{sec:setup}
% =====================================================================
\paragraph{Models and coverage.} We study pretrained \emph{base} code models,
not instruction-tuned chat models: DeepSeek-Coder 1.3B and 6.7B, and StarCoder2
3B. The three models have $24$, $32$, and $30$ layers; cross-model comparisons
are read at matched \emph{relative depth}, never at a shared index. Not every
experiment licenses a claim for every model, because the causal and attribution
instruments have capability and architecture prerequisites
and attribution instruments have capability and architecture prerequisites.
Representation and transfer are reported for all three models. Causal results are
reported for DeepSeek-Coder 6.7B and StarCoder2 3B. Relevance attribution and the
vocabulary readout are reported for DeepSeek-Coder 6.7B. Other configurations
failed a stated gate or fell outside the implemented propagation rules.

\paragraph{Data and ground truth.} Every label is derived from the program's own
structure through a code property graph \citep{yamaguchi2014cpg}: abstract
syntax, a statement-level control-flow graph with exactly resolved join points,
and a data-flow graph whose reaching-definition edges define both binding and the
directed definition--use edge. Labels are exact because the corpus is synthetic
and the graph is known by construction; we validate the extractor against an
independent reaching-definitions implementation (set inclusion in general, exact
equality on straight-line code) and validate every perturbed program by execution
(Appendix~\ref{app:groundtruth}). Graph events are mapped to token indices through
a verified offset table that must reproduce the source exactly; a program whose
round-trip fails is skipped and counted. String matching is not used for
alignment, because a subword vocabulary silently mislabels shadowed names, which
is exactly the phenomenon binding measures.

\paragraph{Read positions.} Hidden states are captured at each decoder block
output and, separately, at the embedding output (layer $-1$), which encodes token
identity before any context mixing and serves as the context-free reference. The
state for a use is read at its last covering token: the first position whose
state can see the whole event under causal attention.

\paragraph{Probe and floors.} The decoding instrument is deliberately weak, 
$\ell_2$-regularized logistic regression ($C{=}0.1$, class-balanced,
standardized features), so that ``decodable'' remains a statement about the
model rather than about probe capacity. Pairwise relations use the endpoint
feature $[\,h_i;\,h_j;\,h_i{-}h_j;\,|h_i{-}h_j|\,]$. All cross-validation is
\code{Stratified\allowbreak GroupK\allowbreak Fold} grouped by source program, so no program has rows in
both train and test; when the corpus is capped, whole groups are dropped. Claims
are made on \emph{selectivity} (accuracy minus a shuffled-label control) and,
critically, against a \emph{model-free surface baseline} that sees only a
$\pm 3$-token window of token ids and the bucketed distance between anchors. This
baseline exists because a first run reproduced $\approx 0.98$ on every task from
local token context alone; the baseline is now a permanent floor, and a
hidden-state result counts only if it beats the baseline on the same stratum.

\paragraph{Negative strata.} Most negative pairs for a relation are separable
from text alone, so negatives are reported in strata from easiest to decisive:
\code{diff\_name} (different identifiers), \code{distance\_matched} (same token
distance), \code{same\allowbreak\_name\allowbreak\_diff\allowbreak\_binding} (same spelling, diff.\ binding),
and \cm{} (two token-identical programs differing by the one binding-flipping
character). Only \cm{} is a clean headline: its members share every fixed-offset
surface cue and one CV group, so by construction no surface feature separates the
labels and the surface probe scores exactly $0.500$ there
(Appendix~\ref{app:strata}).

% =====================================================================
\section{Representation and robustness}
\label{sec:repr}
% =====================================================================

\subsection{Binding is constructed in context (P1)}
On the decisive \cm{} stratum the binding probe is exactly at chance for the
model-free reader and for the embedding layer, then rises sharply through early
blocks and plateaus near $0.984$ in the middle (Table~\ref{tab:binding-layers},
Figure~\ref{fig:strata}). Because the label is absent from the queried token and
from local context, this is a representation claim, not a task-performance claim:
the distinction is \emph{constructed} by contextual processing. The
definition--use edge shows a related profile, accuracy near $0.99$ at layers
7--11 with distance-matched negatives, and $0.96$--$0.99$ for endpoints separated
by $50$--$200$ tokens (Appendix~\ref{app:defuse}), so the model represents both
which definition is active and which use it reaches. Neither probe alone shows
that these relations control the output; that claim is reserved for \S\ref{sec:das}.

\begin{table}[t]
\centering
\small
\caption{Binding-probe accuracy on the \cm{} stratum by depth. The surface
baseline and the embedding layer are pinned at $0.500$ by construction; accuracy
is constructed through early blocks and plateaus in the middle. Peak layers
(L7 for 1.3B; L11--15 for 6.7B) are selected on a calibration split.}
\label{tab:binding-layers}
\begin{tabular}{@{}lcc@{}}
\toprule
Measurement & DeepSeek 1.3B & DeepSeek 6.7B \\
\midrule
Surface baseline (model-free) & \textbf{0.500} & \textbf{0.500} \\
Embedding layer ($-1$)        & \textbf{0.500} & \textbf{0.500} \\
Block 0                        & 0.570 & 0.531 \\
Layer 3                        & 0.961 & 0.914 \\
Peak                           & \textbf{0.984} (L7) & \textbf{0.984} (L11--15) \\
Last layer                     & 0.930 & 0.914 \\
\bottomrule
\end{tabular}
\end{table}

\begin{figure}[t]
\centering
\includegraphics[width=0.66\linewidth]{figs/binding_strata_deepseek-coder-6.7b_core.pdf}
\caption{Binding-probe accuracy by negative stratum and layer
(DeepSeek-Coder 6.7B). Easier strata are partly solvable from surface features;
only \cm{} isolates the semantic distinction, and it sits at $0.500$ at the input.}
\label{fig:strata}
\end{figure}

\subsection{Transfer reflects semantic difficulty (P2)}
A probe trained on clean programs is frozen and re-evaluated on transformed
programs, with ground truth recomputed for every variant; freezing measures
whether the \emph{existing} readout still holds, rather than how learnable the
relation is under each condition (Appendix~\ref{app:freeze}).

\paragraph{Context interference.} Holding the added-token budget approximately
fixed while varying what the inserted text does isolates interference from
length. At $500$ inserted tokens, DeepSeek-Coder 6.7B retains $0.921$ binding
accuracy with inert comments; dead code and lexical decoys give $0.794$--$0.795$,
unrelated competing updates $0.859$, but scope shadows that reuse the tracked
names reduce accuracy to $0.570$. At $1{,}000$ tokens the scope shadow reaches
$0.498$ while all other fillers stay above $0.70$ (Figure~\ref{fig:context}). The
licensed conclusion is comparative: tracked-name scope interference damages
transfer far more than inert distance. It does not show that long context is
universally harmless, nor that a failed frozen probe means the model has
forgotten binding.

\paragraph{Obfuscation ladder.} An execution-verified, cumulative ladder applies
normalization, consistent renaming, opaque predicates, mixed
Boolean--arithmetic rewrites, and control-flow flattening; every variant is
checked observationally equivalent to its base, and the ladder is applied both
individually and composed to attribute failures to steps
(Appendix~\ref{app:obf}). Consistent renaming damages early lexical layers, some
early probes fall below chance, while middle layers survive at $0.85$--$0.90$.
Across models, \emph{control-flow flattening is the largest reproducible failure}
(Table~\ref{tab:obf}, Figure~\ref{fig:obf}). These are frozen-readout results:
they identify where the original representation transfers and thus a confidence
boundary, not a proof that all binding information is erased when the probe fails.

\begin{table}[t]
\centering
\small
\caption{Frozen-probe transfer through the obfuscation ladder (accuracy). Rename
damages the lexical early layers; opaque predicates and arithmetic rewrites add
little; control-flow flattening causes the largest cross-model collapse
(\textbf{bold}).}
\label{tab:obf}
\begin{tabular}{@{}lcccc@{}}
\toprule
Task, model & Rename & Opaque & Arith.\ encoding & Flatten \\
\midrule
Binding, 1.3B        & 0.783 & 0.801 & 0.834 & \textbf{0.555} \\
Binding, 6.7B        & 0.883 & 0.862 & 0.857 & \textbf{0.615} \\
Binding, StarCoder2  & 0.708 & 0.743 & 0.790 & \textbf{0.527} \\
Def--use, 1.3B       & 0.819 & 0.799 & 0.800 & \textbf{0.461} \\
Def--use, 6.7B       & 0.864 & 0.846 & 0.833 & \textbf{0.545} \\
Def--use, StarCoder2 & 0.689 & 0.731 & 0.747 & \textbf{0.402} \\
\bottomrule
\end{tabular}
\end{table}

\begin{figure}[t]
\centering
\begin{subfigure}{0.49\linewidth}
\centering
\includegraphics[width=\linewidth]{figs/context_binding_deepseek-coder-6.7b.pdf}
\caption{Context interference.}
\label{fig:context}
\end{subfigure}\hfill
\begin{subfigure}{0.49\linewidth}
\centering
\includegraphics[width=\linewidth]{figs/obfuscation_levels_deepseek-coder-6.7b.pdf}
\caption{Obfuscation ladder.}
\label{fig:obf}
\end{subfigure}
\caption{Frozen binding-probe transfer (DeepSeek-Coder 6.7B).
(a) Long inert context is far less damaging than scope interference reusing the
tracked names. (b) Middle layers survive renaming; control-flow flattening causes
the largest collapse.}
\label{fig:transfer}
\end{figure}

% =====================================================================
\section{From representation to causal use and attribution}
\label{sec:part3}
% =====================================================================
Sections~\ref{sec:repr} established that binding is represented and measured how
stable that representation is. This section follows the representation into two
distinct consequences on the \emph{same} controlled programs. \S\ref{sec:das}
tests whether replacing one learned binding component makes the answer follow the
installed binding; this is the only experiment that licenses a causal claim.
\S\ref{sec:rlens} leaves the forward computation unchanged and asks whether the
answer score is attributed to the definition the binding selects.
\S\ref{sec:jlens} asks, at the same unprompted state, whether output-vocabulary
margins move with the binding. These three instruments are less familiar than a
linear probe, so we state precisely what each intervenes on, what each reads out,
and what result would falsify its claim.

\subsection{DAS: causal interchange of a binding component (P3)}
\label{sec:das}

\paragraph{What the experiment establishes.} A probe or lens can show that a fact
is present, or present in output coordinates; neither can show that the fact is
\emph{used}, because a representation can be a faithful copy of a computation that
happens elsewhere. Establishing use requires an intervention on the state itself.

\paragraph{The interchange operator.} Given the host state $\hb^{\mathrm{host}}$
(the program the model is actually running) and a donor state
$\hb^{\mathrm{donor}}$ (the counterfactual program), the intervention is
Eq.~\eqref{eq:interchange}: the columns of the orthonormal matrix $\Rmat$ name a
subspace, and the operator installs whatever the donor run has in that subspace
while leaving the orthogonal complement of the host untouched. The design has no
magnitude hyperparameter: unlike an additive edit of chosen norm, an interchange
installs exactly the donor's component, so ``was the dose large enough'' is not a
free question. The realized edit size is instead a \emph{measured} quantity
(\code{edit\_fraction}), reported for every condition. Three algebraic properties
are pinned by unit tests: the component of the host orthogonal to $\mathrm{span}(\Rmat)$
is unchanged; $\hb^{\mathrm{donor}}=\hb^{\mathrm{host}}$ yields exactly the zero
edit, so the no-op control is provably inert; and at $r=d$ the operator returns
the whole donor state, so a whole-state patch is the rank-$d$ limit of the same
operator and the correct ceiling to normalize against. The edit is computed in
float64 and the half-precision hidden state is upcast before the edit and cast
back afterward, so the intervention never occurs in half precision.

\paragraph{How the subspace is learned.} $\Rmat$ is learned in the manner of DAS
\citep{geiger2024das}: only $\Rmat$ is trained, every model parameter is frozen,
and $\Rmat$ is re-orthonormalized inside the computation graph at each step so
that the operator is a true interchange throughout optimization. The state
entering the intervention site is detached, which is exact rather than
approximate, layers before $\ell$ cannot depend on $\Rmat$, so no gradient
should flow to them, and the objective is the negative log-probability of the
token that the donor's value implies, at the final position. Two consequences
follow. A null is strong: it states that \emph{no} $r$-dimensional subspace at
this site behaves as claimed. A positive is weak without controls: because the
method is expressive enough to fit structure that reflects something other than
binding, the controls below are part of the claim, not an appendix to it.

\paragraph{The crossed $2\times2$ that identifies binding.} The design crosses
binding structure with value assignment. Four programs per base are
token-identical except for one character:
\begin{center}
\begin{minipage}{0.46\linewidth}
\begin{lstlisting}
# ARM ab: (outer,inner)=(a,b)
x = a
def f():
    y = b   # target: x = b
    return x
#   -> a (outer) / b (inner)
\end{lstlisting}
\end{minipage}\hfill
\begin{minipage}{0.46\linewidth}
\begin{lstlisting}
# ARM ba: (outer,inner)=(b,a)
x = b
def f():
    y = a   # target: x = a
    return x
#   -> b (outer) / a (inner)
\end{lstlisting}
\end{minipage}
\end{center}
The alignment is fitted on arm \code{ab}, where installing the inner binding
requires the answer to move $a\!\to\!b$, and the claim is read on the held-out arm
\code{ba}, where the \emph{same} binding swap requires the answer to move
$b\!\to\!a$. This crossing is the identifying device: a direction that means
``increase token $b$'' cannot succeed in both arms, whereas a direction carrying
``use the donor's definition'' can. Because the answer \emph{is} the bound value,
there is no arithmetic anywhere and no circularity to forbid; the crossing, not an
arithmetic gap, breaks the confound (Table~\ref{tab:das-identify}).

\begin{table}[t]
\centering
\small
\caption{What each account of the subspace predicts across arms. Only ``which
definition is in scope'' predicts a positive effect in both arms; the observed
DAS pattern matches that row and no other.}
\label{tab:das-identify}
\begin{tabular}{@{}lcc@{}}
\toprule
Account of what the subspace carries & Arm \code{ab} & Arm \code{ba} \\
\midrule
which definition is in scope & positive & \textbf{positive} \\
the token \code{b}           & positive & \textbf{negative} \\
the answer direction         & positive & \textbf{negative} \\
\bottomrule
\end{tabular}
\end{table}

\paragraph{Controls.} Every control is constructed to refute a specific simpler
account (full detail in Appendix~\ref{app:das-controls}). The
\emph{answer-direction} control is decisive: it builds a direction that
explicitly pushes the fitted arm's output from value $a$ toward value $b$, using
vocabulary-lens vectors (because the intervention is partway through the network),
scaled to the DAS edit norm on the same example, and \emph{held fixed} across
arms. It should work in \code{ab} and fail or reverse in \code{ba}. A
\emph{mean-difference} direction (rank-1 span of the mean donor$-$host difference,
no optimizer) is the cheapest non-learned alternative that a learned direction
must dominate, not merely beat. Two random controls separate disruption from
transport: a \emph{dose-matched} random subspace moves the same fraction of
$\|\hb\|$, and a \emph{rank-matched} random subspace is the weaker floor; dose
matching matters because a random rank-$r$ subspace captures only $\approx r/d$ of
the state's energy while a learned one is selected to capture far more, so
rank-matching alone would leave the comparison dose-mismatched in the direction
that manufactures a positive. A no-op verifies the machinery, and a
\emph{def-source} site where the programs are token-identical is a structural zero
where any effect indicates a bug.

\paragraph{Outcome metric.} The primary outcome is \says{}: whether the model's
full-vocabulary argmax is the token the \emph{installed} binding selects, not a
two-way logit margin. When clean accuracy is at ceiling, the model is confident,
and any edit that merely disrupts the state regresses both candidates toward the
middle and \emph{inflates} a two-way margin while transporting nothing; a
disruption cannot systematically produce the correct installed token as the
argmax. The experiment runs behind six gates that refuse downstream stages until
they pass, including a behavioral gate (the model returns the correctly bound
variable, $\geq 0.85$ overall) checked before any instrument is built on top of
it (Appendix~\ref{app:das-gates}).

\paragraph{Result.} DAS reaches $100\%$ installed-answer accuracy in both crossed
arms on the two interpretable architectures (Table~\ref{tab:das}). On
DeepSeek-Coder 6.7B the whole-state ceiling is $\approx 86\%$--$88\%$, so the
rank-1 interchange is above the whole-state patch; the fitted-arm answer-direction
control installs the correct token $27.9\%$ of the time and collapses to $4.3\%$
in the crossed arm, while DAS remains perfect. The mean-difference direction is a
real but weaker baseline that requires a larger edit ($\approx 0.71$ of
$\|\hb\|$ versus $0.48$ for DAS). This is the strongest evidence in the study that
downstream computation \emph{uses} a compact component with the causal role of the
binding variable, at the tested site and layer.

\begin{table}[t]
\centering
\small
\caption{Installed-answer accuracy (\says{}). DAS is fitted on \code{ab} and read
on \code{ba}. Because the correct output token reverses between arms, DAS
succeeding in both cannot be a fixed push toward one token; the answer-direction
control shows exactly the predicted collapse across arms.}
\label{tab:das}
\begin{tabular}{@{}lcccc@{}}
\toprule
Intervention & DeepSeek \code{ab} & DeepSeek \code{ba} & StarCoder2 \code{ab} & StarCoder2 \code{ba} \\
\midrule
\textbf{DAS, rank 1}       & \textbf{100.0\%} & \textbf{100.0\%} & \textbf{100.0\%} & \textbf{100.0\%} \\
Whole-state donor patch    & 85.7\% & 87.9\% & 68.8\% & 67.3\% \\
Mean donor$-$host direction& 76.1\% & 76.8\% & 54.6\% & 54.5\% \\
Answer direction (fixed)   & 27.9\% & 4.3\%  & 44.8\% & 18.4\% \\
Dose-matched random        & 2.1\%  & 1.8\%  & 30.9\% & 30.4\% \\
\bottomrule
\end{tabular}
\end{table}

\subsection{Conservative relevance attribution of the unedited answer (P4)}
\label{sec:rlens}

\paragraph{What the experiment establishes.} DAS shows that changing a rank-1
component at the use site changes the answer; it does not say which source
locations the \emph{unedited} model's answer depends on. The conservative relevance decomposition answers that
second, observational question on the same four-program factorial. It follows DAS
deliberately: if it were the only experiment, a relevance shift could be
overread as weak causal evidence, but here the causal fact is already established
and the conservative relevance decomposition contributes a decomposition of the output score, not a second
intervention.

\paragraph{What it computes.} Fix one output score $s$, the model's score for the
value the binding selects. Ordinary gradients measure local sensitivity but do
not sum to $s$. The conservative relevance decomposition instead modifies only the \emph{backward} computation
with conservation-respecting relevance rules
\citep{bach2015lrp,ali2022xai,achtibat2024attnlrp}: it freezes normalization and
the attention pattern, treats the SiLU nonlinearity as elementwise scaling, and
splits the relevance of a gated MLP equally between its two multiplicative
branches. The forward activations, logits, and emitted token are unchanged. For
the pre-norm, gated-MLP DeepSeek architecture the resulting position relevances
satisfy $\sum_i r_i \approx s$ at every reported layer, so each position's
relevance divided by a \emph{positive} $s$ is a share; shares are summed into
syntactic roles read from the AST (outer definition, inner name, inner value, use
site, signature, \code{return}, suffix, residual). Freezing the attention pattern
is a stated limitation: the method attributes what the fixed pattern transports,
not how queries and keys formed it, so it cannot establish the mechanism
``attention found the correct definition.''

\paragraph{The one-token contrast and its statistic.} Within an arm, the
non-shadowing and shadowing programs differ at exactly one token index out of
$\approx 21$: the inner definition's name. The outer definition, inner value, use
token, signature, and suffix are token-identical and re-verified during the run.
The declared statistic is the \emph{share gained by the token-identical inner
value minus the share retained by the token-identical outer definition}; a
positive value means attribution moved toward the definition entering scope,
measured over text that did not itself change. The instrument is interpreted only
after forward invariance, rule binding, conservation, role partition, and a
same-program zero-redistribution check all pass (Appendix~\ref{app:rlens}). The
rules do not match StarCoder2's normalization and non-gated MLP, so the pipeline
refuses to report relevance semantics there; and conservation alone is insufficient
when shares are reported, so the sign of $s$ is checked separately, which fails
often enough on the 1.3B model that its shares are not interpreted.

\paragraph{Result.} On DeepSeek-Coder 6.7B the shift has the predicted sign on all
$280$ held-out bases. At the first measured layer the newly active inner value
gains $\approx 4.9\%$ of the answer score and the newly inactive outer definition
loses $\approx 7.8\%$, a combined redistribution of $\approx 12.6\%$; the shift
rises to $\approx 21.9\%$ in the middle and declines to $\approx 2.5\%$ near the
output. The single changed name token carries only $\approx 1.5\%$ of the
movement. The effect appears in both value arms, survives scoring both members at
one fixed output token, reverses when the competing value is scored, and is
absent in same-binding controls. It is therefore consistent with attribution
following the active definition rather than responding locally to the changed
name. This is observational evidence that complements, but does not amplify, the
causal status of DAS; because all bases share one template, the effect is best
read as a property of the template-level contrast (Appendix~\ref{app:rlens}).

\subsection{Archived prompt-specific cotangent readout of lexical alignment (P5)}
\label{sec:jlens}

\paragraph{What the experiment establishes.} The final question is whether the
binding state is expressed in the model's \emph{output vocabulary} at the
unprompted use token, with no appended question. This archived instrument is a
prompt-specific cotangent vocabulary readout, not the published J-lens. The logit lens
reads a hidden state directly through the unembedding
\citep{nostalgebraist2020logitlens}, the prompt-specific cotangent readout estimates how the remaining layers
map a change at the read position to the output, which is the relevant map when
the state is read mid-network. We report the plain logit lens alongside it to
separate vocabulary alignment from an effect introduced by the Jacobian
correction.

\paragraph{Design.} For each of nine predeclared single-token pairs, four scope
(\code{local/global}, \code{nested/module}, \code{inner/outer}, and one
additional), three positional, and two action contrasts, we measure, per layer
and per value arm, the fraction of held-out programs on which the
inner-minus-outer margin moves in the predicted direction when only the binding
changes. The two value arms are the principal control: activating the inner
definition returns literal \code{b} in \code{ab} but literal \code{a} in
\code{ba}, so a margin that shifts the same binding-relative way in both arms is
not following one answer token. Pairs are grouped into families so scope
contrasts can be compared against positional and action contrasts. The positive
control is a calibration-fitted binding probe on the same held-out states, which
confirms the read position contains binding information; every word pair must
preserve its declared tokenization (Appendix~\ref{app:jlens}).

\paragraph{Result and its limit.} The instrument is mechanically valid: the
matched binding probe is $1.000$ at L8, L12, L16, L20, and L24, and several scope
pairs follow the binding strongly in both arms, \code{nested/module} reaches
$1.000/1.000$ at L16 and \code{local/global} $0.996/1.000$ at L20, persisting past
$0.98$ at L24 (Table~\ref{tab:jlens}). The crossing makes these associations
meaningful. They do not, however, isolate scope semantics. Under the single
template the active definition simultaneously becomes later, nearer, second, and
retained, so the positional \code{later/earlier} and action \code{replaced/kept}
contrasts move as strongly as the scope pairs while the literal \code{inner/outer}
pair is weak. The supported conclusion is therefore
\emph{binding-associated lexical alignment}, not that the representation
verbalizes ``scope.'' A plausible-sounding rationale that uses correct semantic
vocabulary is thus not, on this evidence, a faithful readout of the causal state.

\begin{table}[t]
\centering
\small
\caption{Archived prompt-specific cotangent crossed-arm reversal rate (fraction of held-out programs on which
the margin moves with the binding), DeepSeek-Coder 6.7B, at representative
layers. Scope pairs track the binding, but a positional confound tracks it as
strongly, so alignment is not uniquely scope.}
\label{tab:jlens}
\begin{tabular}{@{}llcc@{}}
\toprule
Pair, layer & Family & Reversal \code{ab} & Reversal \code{ba} \\
\midrule
\code{nested/module}, L16   & scope      & 1.000 & 1.000 \\
\code{local/global}, L20    & scope      & 0.996 & 1.000 \\
\code{nested/module}, L24   & scope      & 0.993 & 0.996 \\
\code{local/global}, L24    & scope      & 0.986 & 0.993 \\
\code{later/earlier}, L16   & positional & 1.000 & 1.000 \\
\bottomrule
\end{tabular}
\end{table}

% =====================================================================
\section{Discussion}
\label{sec:discussion}
% =====================================================================
No single result establishes semantic grounding; their conjunction is what is
informative. The exact embedding and surface nulls show that binding is not
locally encoded in the queried token. The rise through layers shows contextual
construction. Renaming resistance in middle layers shows partial independence from
identifier strings. Greater sensitivity to competing scope than to inert distance
ties failure to semantic difficulty rather than to length. The crossed DAS
intervention shows causal use that is separable from a fixed answer token, and the
the relevance decomposition links the unedited answer to the active definition. Together these are
difficult to reconcile with a purely lexical account.

The same results also reject a strong symbolic account. The frozen representation
is not invariant to all meaning-preserving transformations: flattening and scope
interference sharply reduce transfer, early layers can be misled by renaming, and
the vocabulary alignment is confounded across scope, position, and replacement.
The models compute a useful binding relation, but that computation remains
entangled with learned lexical and structural conventions. The practical reading
is a \emph{conditional} one. Tasks dominated by ordinary lexical patterns can hide
semantic weaknesses, because those patterns usually agree with program meaning;
confidence should fall on inputs where lexical and semantic evidence diverge, 
adversarial naming, unfamiliar control-flow lowering, deeply competing scopes, 
which are exactly the regimes in which the measured readout loses transfer. An
assurance case for a code model should therefore test representation, invariance,
causal use, and failure boundaries separately, and restrict or escalate actions
outside the validated regime, rather than infer trust from one correct output or
one plausible explanation.

% =====================================================================
\section{Conclusion}
% =====================================================================
We tested semantic grounding with a binding construction whose surface floor is
exactly $0.500$. The models form a linearly recoverable binding relation and use
a rank-1 component of it when selecting an answer. DeepSeek-Coder 6.7B also
assigns answer relevance toward the active definition, and its unprompted state
aligns with several lexical contrasts after answer identities are reversed.
Competing scope and flattened control flow reduce transfer. The lexical readout
does not isolate scope from position. These results support trust within the
tested conditions and identify concrete limits on that trust.

% =====================================================================
\section*{Limitations and threats to validity}
\addcontentsline{toc}{section}{Limitations and threats to validity}
% =====================================================================

\paragraph{External validity.} The causal experiment is a small synthetic Python
construction with one lexical mutation, one use site, and scalar-literal answers.
That control is what makes the causal claim possible, but it does not cover
repository scale, interprocedural aliasing, objects, closures, mutation through
containers, exceptions, concurrency, or dynamic dispatch. The next step is to
embed \cm{} counterfactuals in diverse real code instead of adding generated names.
A general form-preserving, meaning-breaking construction, a program that presents
as one thing and computes another, does not exist in this corpus, and it is the
cell closest to the adversarial motivation.

\paragraph{Model and instrument scope.} Representation and robustness span three
model families where reported, but the conservative relevance attribution claim holds only for
DeepSeek-Coder 6.7B (the 1.3B model has non-positive selected scores on part of
the condition; StarCoder2 is outside the implemented relevance rules), and the
crossed-arm vocabulary-lens result is currently 6.7B only. Cross-model numbers are
read at matched relative depth, but depth alignment is itself an approximation.

\paragraph{What a probe failure does and does not mean.} Frozen-probe transfer
failure shows that a clean-trained linear readout no longer transfers; it does not
prove the information was erased, since a refitted or nonlinear probe might recover
it. Symmetrically, high probe accuracy does not prove use, which is why DAS is
necessary. DAS is itself local: it identifies a successful rank-1 interchange at
one site and layer, not a unique global mechanism, and its superiority to a
whole-state patch suggests the full state carries opposing components, an
explanation that remains untested.

\paragraph{Attribution and lens caveats.} The conservative relevance decomposition is rule-dependent and
observational; conservation verifies an accounting identity, not that the
resulting partition is the unique causal explanation, and freezing the attention
pattern precludes attributing the query/key formation. Because every base shares
one template, a mismatched-base control reproduces the mean shift, so effect sizes
should be read as properties of the controlled contrast. The archived prompt-specific cotangent readout has limited
early-layer stability (independent-build sign agreement exceeds $0.90$ only at
L24), its next-token validation is small, and its single template changes scope,
order, distance, and replacement together; broader templates that break these
correlations are required before naming the alignment as scope.

\paragraph{Grounding is necessary, not sufficient, for safety.} A model can
represent binding correctly while following a malicious instruction, calling an
unsafe tool, or optimizing the wrong objective. Access control, sandboxing,
approval boundaries, and conventional software assurance remain necessary; this
paper contributes evidence about one component of trustworthy reasoning, not a
safety case.

% References
\bibliography{references_revised}

\clearpage
\appendix
% ============================ APPENDIX ============================
\section{Ground truth: code property graph and cross-checks}
\label{app:groundtruth}
Every label is derived from the program's own structure rather than from a
learned labeller, for two reasons. First, a probe trained on approximate labels
measures the approximation as much as the model, so the primary corpora are
synthetic programs whose graph is known by construction; real-code transfer is
reported only as an upper bound. Second, the exact $0.500$ floor exists only
because the generator can emit two programs with identical surface and different
graphs, which requires authoring the graph.

Three independent cross-checks guard the extractors. (i) Reaching definitions are
compared against \code{beniget}, an independently written analysis. The two answer
slightly different questions---ours resolves each use to the single most recent
reaching definition, \code{beniget} returns all possibly-reaching definitions
across branches---so the sound comparison is set inclusion (our edges $\subseteq$
\code{beniget}'s) with exact equality on straight-line code. This caught a real
bug: uses in self-referential updates such as \code{b = b + a} were linked to the
same-line target rather than the prior definition; the extractor now resolves in
execution order. (ii) Every obfuscation variant is executed and checked
observationally equivalent to its base. (iii) The binding factorials are decided
twice, by execution and by a scope-aware reference interpreter, and both must
agree with the intended label.

\paragraph{Tokenizer integrity.} \code{AutoTokenizer} on \code{transformers} 5.x
can silently resolve DeepSeek-Coder to a slow path that mis-tokenizes code
(\code{def func} $\to$ \code{['de','ff','unc']}), producing corrupt activations
without raising an error. The loader therefore uses \code{PreTrainedTokenizerFast}
and rejects any tokenizer that fails an exact code round-trip; results predating
this guard were rerun.

\section{Excluded relation: control dependence}
\label{app:controldep}
Control dependence has a \emph{measured} surface floor of $0.927$: a statement's
guard is usually its nearest enclosing \code{if}, so a $\pm 3$-token window plus
indentation recovers most of the relation with no model. The models decode it at
ceiling (AUC $0.999$), but because the floor is already near ceiling, this
constrains nothing about semantic representation. Under the operational definition
(\S\ref{sec:formulation}) the relation is mostly syntactic, and it is therefore
not reported as a result. This is the definition doing work: a criterion that
never excludes anything is not a criterion.

\section{Probes, selectivity, and negative strata}
\label{app:strata}
The probe is $\ell_2$-regularized logistic regression ($C{=}0.1$, class-balanced,
\code{saga}, \code{max\_iter}$=2000$, \code{tol}$=10^{-3}$) on standardized
features; convergence is recorded per row and the stage fails if any reported fit
did not converge. A high-capacity probe is avoided on purpose: it could learn the
relation itself, converting a statement about the model into a statement about the
probe.

\paragraph{Selectivity.} For every probe the identical classifier is retrained on
shuffled labels and \code{selectivity = accuracy $-$ control\_accuracy} is
reported; claims are made on selectivity. Labels are shuffled within program for
pairwise and per-token tasks (preserving each program's label mix and destroying
only the token$\to$label pairing) and across programs for example-level tasks.

\paragraph{Surface baseline.} A model-free reader sees only the $\pm 3$-token
window of token ids around each anchor plus the bucketed anchor distance, fit with
the same grouped CV, reported per stratum. It exists because a first run scored
$\approx 0.98$ on every task and layer, which the baseline reproduced, proving the
templated corpus was leaking labels through local token context.

\paragraph{Strata.} Table~\ref{tab:strata} defines the negative strata and what
each controls. Only \cm{} is a clean headline: both members are token-identical
except the one binding-flipping character, the anchor windows and token distance
are identical, and both share one CV group, so by construction no surface feature
separates the labels and the surface probe scores exactly $0.500$.

\begin{table}[h]
\centering\small
\caption{Negative strata for the binding probe, easiest to decisive.}
\label{tab:strata}
\begin{tabular}{@{}lll@{}}
\toprule
Stratum & Construction & Controls for \\
\midrule
\code{diff\_name} & different variable names ($\leq 3\times$ positives) & trivial baseline \\
\code{distance\_matched} & negatives at the positives' token distance & positional shortcuts \\
\code{same\_name\_diff\_binding} & same spelling, different actual binding & the name-identity shortcut \\
\textbf{\cm{}} & token-identical pair, one binding-flipping character & every fixed-offset surface cue \\
\bottomrule
\end{tabular}
\end{table}

\section{Frozen transfer and the obfuscation ladder}
\label{app:freeze}\label{app:obf}
Robustness probes trained on clean base programs are evaluated unchanged on
transformed programs; they are never retrained per condition. Retraining would
measure how learnable the relation is under each condition---a different, easier
question that lets the probe find a new per-condition shortcut and conceals the
representational failure being sought. Ground truth is recomputed for every
variant, since some transformations (competing updates; inserted opaque branches;
flattening) genuinely change the graph.

The ladder is cumulative and Tigress-inspired, implemented natively for Python
(Table~\ref{tab:ladder}). Every variant is executed and checked observationally
equivalent to its base, and all levels of a base are kept or dropped together so
per-level curves compare an identical set of programs. Because a cumulative ladder
cannot attribute a failure (its last rung contains every earlier one), the same
four rewrites are also applied individually; each reported cell yields
\code{delta\_clean} (cost vs.\ clean held-out), \code{delta\_previous} (marginal
cost of the added step), and \code{delta\_atomic} (cumulative minus atomic, the
interaction). The interaction is read against a measured draw-noise floor: two
conditions apply the identical transformation under independent random draws, and
their difference is that floor. Each variant's transformations are read off its
own AST and must equal exactly what its condition declares.

\begin{table}[h]
\centering\small
\caption{The obfuscation ladder. Each variant is execution-verified equivalent to
its base.}
\label{tab:ladder}
\begin{tabular}{@{}cll@{}}
\toprule
Level & Transformation & What it changes \\
\midrule
0 & normalize (\code{ast.unparse}) & formatting only (baseline) \\
1 & + consistent renaming of every local & identifier strings \\
2 & + dead branches under false opaque predicates & apparent control structure \\
3 & + mixed Boolean--arithmetic rewriting & expression form \\
4 & + control-flow flattening into a dispatch loop & the entire control skeleton \\
\bottomrule
\end{tabular}
\end{table}

\section{Definition--use edge results}
\label{app:defuse}
The definition--use experiment labels whether a directed reaching-definition edge
connects one token to another, with distance-matched negatives so a nearest-token
heuristic cannot solve it. Accuracy peaks near $0.99$ at layers 7--11 and remains
$\approx 0.96$--$0.99$ for endpoints separated by $50$--$200$ tokens
(Figure~\ref{fig:defuse}). With binding, this gives two related structural
observations: the model represents which definition is active and which use it
reaches.

\begin{figure}[h]
\centering
\includegraphics[width=0.6\linewidth]{figs/defuse_distance_deepseek-coder-6.7b_core.pdf}
\caption{Definition--use probe accuracy by endpoint distance
(DeepSeek-Coder 6.7B, \cm{}-style distance-matched negatives).}
\label{fig:defuse}
\end{figure}

\section{DAS: controls}
\label{app:das-controls}
Table~\ref{tab:das-controls} lists every control and the specific account it
refutes. Two points deserve emphasis. The \emph{answer-direction} control uses
J-lens vectors rather than raw final-layer unembedding rows because the
intervention occurs partway through the network and the J-lens estimates how the
remaining layers map a change there to the output; a synthetic donor makes the
interchange an exact push along this direction, and its length is set equal to the
DAS edit on the same example, so DAS is not given a larger dose. The
\emph{concentration} diagnostic reports the share of a basis's mass on its largest
residual-stream dimensions: transformer residual streams contain a few
massive-activation dimensions, and an unconstrained low-rank fit can align with one
to produce a large logit shift while transporting nothing; a basis spread over the
stream gives $\approx \text{top-}k/d$, one riding a rogue dimension approaches
$1.0$.

\begin{table}[h]
\centering\small
\caption{DAS controls and the account each refutes.}
\label{tab:das-controls}
\begin{tabular}{@{}lp{0.62\linewidth}@{}}
\toprule
Control & What it refutes \\
\midrule
\code{whole\_state} & the rank-$d$ limit; the per-arm ceiling. Its being alive in both arms is what makes a null in either arm interpretable. \\
\code{mean\_difference} & the cheapest non-learned rank-1 direction; a learned direction must dominate it, not merely beat zero. \\
\code{answer\_direction} & ``the subspace just pushes toward the fitted arm's answer token''; must work in \code{ab} and fail/reverse in \code{ba}. \\
\code{random\_norm} & disruption; a random subspace moving the same fraction of $\|\hb\|$. \\
\code{random\_rank} & the weaker, rank-matched floor. \\
\code{noop} & machinery; provably exactly $0.00\mathrm{e}{+}00$. \\
\code{def\_source} site & a structural zero at a token-identical site; any effect is a bug. \\
\bottomrule
\end{tabular}
\end{table}

\section{DAS: gates}
\label{app:das-gates}
Each gate refuses to run downstream stages until it passes; thresholds were fixed
before any large-model run. H1 exists because of a recorded lesson: verify the
model can do the task before building an instrument on it. On DeepSeek-Coder 6.7B
all gates pass: H0 (execution and the reference interpreter agree, and every
invariant including the arm crossing holds) at $1.0000$ of $400$ bases; H1
behavioral accuracy $1.000\,[1.000,1.000]$; H2 binding decodable at $1.000$
(selectivity $0.524$) at L8 against a measured surface baseline of $0.500$; H3
whole-state flip margins $+4.78/+4.80$ with flip rates $0.86/0.88$ in the two
arms; H4 rank-1 interchange at $189\%$ of the whole-state ceiling on the fitted
arm; H5 transfer to the held-out arm at $100\%$ installed-answer with the
answer-direction control failing there.

\begin{table}[h]
\centering\small
\caption{DAS gates and thresholds.}
\label{tab:das-gate-defs}
\begin{tabular}{@{}clc@{}}
\toprule
Gate & Asserts & Threshold \\
\midrule
H0 & execution and reference interpreter agree; all invariants incl.\ arm crossing & $\geq 0.999$ of bases \\
H1 & the model returns the correctly bound variable & $\geq 0.85$ overall, $\geq 0.75$ per cell \\
H2 & binding decodable at the use anchor & $\geq 0.80$, $\geq 0.10$ over surface baseline \\
H3 & whole-state interchange flips the answer in both arms & CI $>0$, flip rate $\geq 0.25$ \\
H4 & low-rank interchange beats controls on the training arm & $\geq 50\%$ of ceiling \\
H5 & the same subspace transfers to the held-out arm, where \code{answer\_direction} fails & $\geq 50\%$ of ceiling \\
\bottomrule
\end{tabular}
\end{table}

\section{R-lens: checks and controls}
\label{app:rlens}
The R-lens is interpreted only after five checks pass: forward invariance
(installing the backward rules changes no forward output), rule binding (the
normalization, attention, and gated-MLP rules attach to the intended modules),
conservation (relevance sums to the selected score at every reported layer), role
partition (every encoded token belongs to exactly one syntactic role), and a
same-program reread that must produce exactly zero redistribution. Conservation
holds with median $|\rho-1|\approx 10^{-7}$ across reported layers on the tested
DeepSeek architectures; disabling any rule (no-LN, no-attention, no-fp32) breaks
conservation, confirming the rules are load-bearing. Table~\ref{tab:rlens-controls}
lists the substantive controls. The mismatched-base control has limited power
here: every base shares one template, so mismatching still compares non-shadowing
with shadowing, which bounds the finding to a population-level template contrast
rather than showing it spurious.

\begin{table}[h]
\centering\small
\caption{R-lens controls and the mechanism each isolates.}
\label{tab:rlens-controls}
\begin{tabular}{@{}lp{0.6\linewidth}@{}}
\toprule
Control & Mechanism isolated \\
\midrule
token-identical statistic & excludes direct relevance at the changed name, length, and positional drift \\
crossed \code{ab}/\code{ba} arms & the scored token moves oppositely; agreement rules out a fixed output token \\
\code{fixed\_a}/\code{fixed\_b} & both programs scored at one literal token id; removes output-token identity \\
competing target & scoring the non-selected value should reverse the shift \\
same-binding contrasts & values change, binding does not; should stay flat \\
random orientation & randomly reversing pair direction should destroy the signed mean \\
mismatched bases & tests dependence on the exact pairing (limited power here) \\
same-program reread & structural zero \\
\bottomrule
\end{tabular}
\end{table}

\section{J-lens: validation}
\label{app:jlens}
The vocabulary-lens run passes its mechanical gates: all exactness cells pass, all
nine word pairs preserve their declared tokenization, and the matched binding
probe is $1.000$ at the reported layers, confirming that the read state contains
binding information. The plain logit lens is reported alongside the
Jacobian-corrected lens so that vocabulary alignment can be separated from an
effect introduced by the correction. Independent-build sign agreement exceeds
$0.90$ only at L24, and the held-out lexical next-token validation contains few
positions, so early-layer margins are treated as unstable. Because the single
template changes scope, textual order, distance, and replacement status together,
even a high crossed-arm rate supports binding-associated alignment only, not
uniquely scope-semantic or faithful verbalization; the positional and action
families are reported precisely so this confound is visible rather than hidden.

\section{Statistics and reproducibility}
\label{app:stats}
Uncertainty is a cluster bootstrap over source programs, never over rows, because
rows from one program share hidden vectors and a row-level bootstrap gives
intervals that are too narrow in the direction that makes a null look like a
finding; control comparisons are paired on the same rows so the difference carries
the interval. For heavy-tailed quantities the median and sign are reported beside
the mean. Layer and site are chosen on a calibration split and recorded before any
test number is read; frozen artifacts (probes, lenses, subspaces, discovered token
sets) are written to disk on the calibration side and read back on the evaluation
side, so the separation is a filesystem boundary. Weak gates exit non-zero on a
failed check; hard gates refuse to run (exit 2) unless their prerequisites have
passed, and any \code{--override-gate} is recorded permanently in the manifest and
in every output row. Seed $42$ is used throughout (generator, CV splits,
subsampling, bootstrap); every stage writes a manifest with the git revision,
arguments, and wall-clock time; and all figures and tables regenerate from the
tidy CSVs alone, so the chain from raw data to published figure is auditable end
to end. A passing gate certifies that a measurement is well-formed at that step;
it is not itself a scientific claim, and a valid instrument must still pass when no
semantic effect is present.


\end{document}
